\subsection{PAC-Bayes and Information-Theoretic Bounds}

PAC-Bayes bounds analyze randomized predictors by comparing a data-dependent posterior distribution on hypotheses to a fixed prior.

\begin{definition} \label{def:sl-randomized-predictor} \hlt{Randomized Predictor}\\
Let $\mathcal{H}$ be a hypothesis space.  A distribution $Q$ on $\mathcal{H}$ defines a randomized predictor by sampling $h\sim Q$ and predicting with $h$.  Its risk is
\[
R(Q)=\mathbb{E}_{h\sim Q}R(h).
\]
\end{definition}

\begin{theorem} \label{thm:sl-pac-bayes-kl-form} \hlt{PAC-Bayes Bound Form}\\
Let $\Pi$ be a prior on $\mathcal{H}$ independent of the sample.  For bounded losses, with probability at least $1-\delta$, every posterior $Q$ satisfies a bound of the form
\[
R(Q)\leq \widehat{R}_S(Q)
+O\!\left(
\sqrt{\frac{\operatorname{KL}(Q\Vert \Pi)+\log(1/\delta)}{n}}
\right).
\]
\end{theorem}

\begin{remark} \label{rem:sl-pac-bayes-prior-posterior} \hlt{Prior and Posterior}\\
The prior in PAC-Bayes is not required to be a subjective Bayesian prior.  Its role is to anchor a complexity penalty: posteriors close to the prior are simpler in the sense of relative entropy.
\end{remark}

\begin{definition} \label{def:sl-information-generalization-gap} \hlt{Information-Theoretic Generalization}\\
If $W=A(S)$ is the learned output of an algorithm, information-theoretic bounds control the generalization gap by quantities such as the mutual information $I(W;S)$.  Such bounds measure how much the trained model reveals about the particular sample.
\end{definition}
